% ============================================================
%  HEADER / PREAMBLE — NaT1_PAB
%  Compile with: pdflatex (2-3 passes, for TOC/refs/bookmarks)
% ============================================================
\documentclass[11pt,twoside,a4paper]{article}

% ============================================================
% DOCUMENT GEOMETRY
% ============================================================
\usepackage[a4paper,left=1cm,right=1cm,top=1.6cm,bottom=1.6cm,twoside,headheight=15pt]{geometry}

% ============================================================
% ENCODING & LANGUAGE
% ============================================================
\usepackage[utf8]{inputenc}
\usepackage[german]{babel}   % new German orthography ("neue deutsche Rechtschreibung")

% ============================================================
% MATH PACKAGES & FONTS
% ------------------------------------------------------------
% IMPORTANT LOAD ORDER: amsmath/amssymb and any package that
% defines math operators/symbols (trfsigns, trsym) MUST be
% loaded BEFORE newtxmath. newtxmath redefines many low-level
% math commands and will throw "already defined" errors if it
% is loaded first (classic conflicts: \Bbbk, \openbox, \laplace).
% ============================================================
\usepackage{amsmath}          % core math environments/commands
\usepackage{amssymb}          % extra math symbols

\usepackage{libertine}             % elegant serif font (Linux Libertine clone)
\usepackage[libertine]{newtxmath}  % math fonts matched to Libertine
\usepackage{tgheros}               % modern sans-serif (TeX Gyre Heros)

% trfsigns and newtxmath both try to define \laplace without
% checking if it already exists -> free the name so newtxmath
% can redefine it without a fatal "already defined" error.
\makeatletter
\let\laplace\relax
\makeatother

\usepackage{trfsigns, trsym}  % Laplace/Fourier transform symbols

% ============================================================
% GRAPHICS, TABLES & LAYOUT HELPERS
% ============================================================
\usepackage{graphicx}     % include images
\usepackage{makecell}     % helper for building tables
\usepackage{bigints}      % large integral signs
\usepackage{fancyhdr}     % custom headers and footers
\usepackage{pgf}
\usepackage{tabularx}     % flexible-width table columns
\usepackage{multirow}     % merge table cells across rows
\usepackage{efbox}        % nice colored boxes
\usepackage{pdflscape}    % landscape pages that rotate in PDF viewers
\usepackage{rotating}     % rotate arbitrary content
\usepackage{multicol}     % multi-column layouts
\usepackage{tikz}         % vector graphics / diagrams
\usepackage{pgfplots}
\pgfplotsset{compat=1.18}
\usepackage{pdfpages}     % include external PDF pages
\usepackage{float}        % better control over figure/table placement

% ------------------------------------------------------------
% PFG Settings
% ------------------------------------------------------------
\pgfplotsset{
  spektrum/.style={
    width=\linewidth, height=4.5cm,
    axis lines=middle,
    axis line style={-latex, thick},
    xmin=-10.8, xmax=10.8,
    xtick={-10,-8,...,10},              % Beschriftung nur gerade k
    minor xtick={-9,-7,...,9},          % ungerade k nur als Strich
    tick label style={font=\scriptsize},
    every axis x label/.style={at={(current axis.right of origin)}, anchor=west},
    every axis y label/.style={at={(current axis.above origin)}, anchor=south},
    xlabel={$\omega$},
    clip=false,
  },
  stab/.style={ycomb, red, very thick, mark=*, mark options={fill=red, scale=0.8}},
  huelle/.style={blue, dashed, thin, samples=80},
}

% ------------------------------------------------------------
% SI units (siunitx)
% ------------------------------------------------------------
\usepackage{siunitx}
\sisetup{
	range-phrase = --,   % separator for value ranges, e.g. 1--5
	range-units  = single
}

% ============================================================
% COLOR PALETTE (template style)
% ============================================================
\usepackage{xcolor}
\definecolor{primary}{HTML}{1B3A6B}   % deep blue
\definecolor{accent}{HTML}{2FA4A9}    % teal
\definecolor{accent2}{HTML}{E8813A}   % warm orange
\definecolor{lightgray}{HTML}{F4F5F7}
\definecolor{darkgray}{HTML}{333333}

% ============================================================
% HYPERLINKS & PDF METADATA
% ============================================================
\usepackage[hidelinks]{hyperref}
\hypersetup{
	pdfauthor  = {\authorinfo},
	pdftitle   = {\subjectinfo},
	colorlinks = true,
	linkcolor  = primary,
	urlcolor   = accent,
	citecolor  = accent2
}
\author{\authorinfo}
\title{\subjectinfo}

% ============================================================
% HEADER & FOOTER (fancyhdr)
% ============================================================
\pagestyle{fancy}
\fancyhf{}                              % clear all header/footer fields
\renewcommand{\headrulewidth}{0.6pt}    % rule under the header
\renewcommand{\footrulewidth}{0pt}      % no rule above the footer

\fancyhead[C]{\small\color{darkgray}\nouppercase\leftmark}  % current section name

\fancyfoot[C]{\small\thepage}
\fancyfoot[L]{\footnotesize\color{darkgray}{\authorinfo}}
\fancyfoot[R]{\footnotesize\color{darkgray}\subjectinfo}

% ============================================================
% PARAGRAPH STYLE
% ============================================================
% Disable paragraph indentation (international style, most
% lecturers write this way too).
\setlength{\parindent}{0pt}

% ============================================================
% CUSTOM MATH OPERATORS
% ============================================================
\DeclareMathOperator{\sinc}{sinc}
\DeclareMathOperator{\sgn}{sgn}
\DeclareMathOperator{\Real}{Re}
\DeclareMathOperator{\Imag}{Im}
%\DeclareMathOperator{\e}{e}
\DeclareMathOperator{\cov}{cov}
\DeclareMathOperator{\PolyGrad}{PolyGrad}

% Macro for the upright 'd' used in integrals/differential equations
\newcommand*{\diff}{\mathop{}\!\mathrm{d}}

% ============================================================
% BOXED EQUATIONS
% ============================================================
\usepackage{empheq}
\newcommand{\boxedeq}[1]{\begin{empheq}[box={\fboxsep=6pt\fbox}]{align*}#1\end{empheq}}

% ============================================================
% TABLE SETTINGS
% ============================================================
% Default row height/stretch factor for tables
\newcommand{\arraystretchOriginal}{1}
\renewcommand{\arraystretch}{\arraystretchOriginal}

% Centered column type for tabularx (usage: Y instead of X)
\newcolumntype{Y}{>{\centering\arraybackslash}X}

% ============================================================
% SECTION TITLES (template style: colored titles + rule under \section)
% ============================================================
\usepackage{titlesec}

\setcounter{secnumdepth}{4}  % number headings down to \paragraph
\setcounter{tocdepth}{2}     % but only list sections/subsections (x.x) in the TOC

\titleformat{\section}
	{\normalfont\Large\bfseries\color{primary}}
	{\thesection}{0.6em}{}[{\color{accent}\titlerule[1.2pt]}]
\titleformat{\subsection}
	{\normalfont\large\bfseries\color{accent}}
	{\thesubsection}{0.5em}{}
\titleformat{\subsubsection}
	{\normalfont\bfseries\color{darkgray}}
	{\thesubsubsection}{0.5em}{}
\titleformat{\paragraph}
	{\normalfont\normalsize\bfseries\color{accent2}}
	{\theparagraph}{1em}{}

\titlespacing*{\section}{0pt}{1.4em}{0.8em}
%\titlespacing*{\paragraph}{0pt}{3.25ex plus 1ex minus .2ex}{1.5ex plus .2ex}

% Reduce the vertical space LaTeX normally inserts before
% subsections/subsubsections, for a more compact layout.
%\let\OriginalSection=\section
%\renewcommand{\section}[1]{\OriginalSection{#1}\vspace{-0.0em}}

\let\OriginalSubSection=\subsection
\renewcommand{\subsection}[1]{\OriginalSubSection{#1}\vspace*{-0.8em}}

\let\OriginalSubSubSection=\subsubsection
\renewcommand{\subsubsection}[1]{\OriginalSubSubSection{#1}\vspace*{-0.8em}}

% ============================================================
% TABLE OF CONTENTS
% ============================================================
% Redefine \tableofcontents to just dump the .toc file, without
% printing an extra "Contents" heading or starting a new page
% (useful when the TOC title/placement is handled manually
% elsewhere in the document).
\makeatletter
\renewcommand\tableofcontents{%
	\@starttoc{toc}%
}
\makeatother